% Options for packages loaded elsewhere
\PassOptionsToPackage{unicode}{hyperref}
\PassOptionsToPackage{hyphens}{url}
\documentclass[
]{article}
\usepackage{xcolor}
\usepackage{amsmath,amssymb}
\setcounter{secnumdepth}{-\maxdimen} % remove section numbering
\usepackage{iftex}
\ifPDFTeX
  \usepackage[T1]{fontenc}
  \usepackage[utf8]{inputenc}
  \usepackage{textcomp} % provide euro and other symbols
\else % if luatex or xetex
  \usepackage{unicode-math} % this also loads fontspec
  \defaultfontfeatures{Scale=MatchLowercase}
  \defaultfontfeatures[\rmfamily]{Ligatures=TeX,Scale=1}
\fi
\usepackage{lmodern}
\ifPDFTeX\else
  % xetex/luatex font selection
\fi
% Use upquote if available, for straight quotes in verbatim environments
\IfFileExists{upquote.sty}{\usepackage{upquote}}{}
\IfFileExists{microtype.sty}{% use microtype if available
  \usepackage[]{microtype}
  \UseMicrotypeSet[protrusion]{basicmath} % disable protrusion for tt fonts
}{}
\makeatletter
\@ifundefined{KOMAClassName}{% if non-KOMA class
  \IfFileExists{parskip.sty}{%
    \usepackage{parskip}
  }{% else
    \setlength{\parindent}{0pt}
    \setlength{\parskip}{6pt plus 2pt minus 1pt}}
}{% if KOMA class
  \KOMAoptions{parskip=half}}
\makeatother
\usepackage{longtable,booktabs,array}
\usepackage{caption}
\captionsetup[table]{skip=6pt}
\newcounter{none} % for unnumbered tables
\usepackage{calc} % for calculating minipage widths
% Correct order of tables after \paragraph or \subparagraph
\usepackage{etoolbox}
\makeatletter
\patchcmd\longtable{\par}{\if@noskipsec\mbox{}\fi\par}{}{}
\makeatother
% Allow footnotes in longtable head/foot
\IfFileExists{footnotehyper.sty}{\usepackage{footnotehyper}}{\usepackage{footnote}}
\makesavenoteenv{longtable}
\usepackage{graphicx}
\makeatletter
\newsavebox\pandoc@box
\newcommand*\pandocbounded[1]{% scales image to fit in text height/width
  \sbox\pandoc@box{#1}%
  \Gscale@div\@tempa{\textheight}{\dimexpr\ht\pandoc@box+\dp\pandoc@box\relax}%
  \Gscale@div\@tempb{\linewidth}{\wd\pandoc@box}%
  \ifdim\@tempb\p@<\@tempa\p@\let\@tempa\@tempb\fi% select the smaller of both
  \ifdim\@tempa\p@<\p@\scalebox{\@tempa}{\usebox\pandoc@box}%
  \else\usebox{\pandoc@box}%
  \fi%
}
% Set default figure placement to htbp
\def\fps@figure{htbp}
\makeatother
\setlength{\emergencystretch}{3em} % prevent overfull lines
\providecommand{\tightlist}{%
  \setlength{\itemsep}{0pt}\setlength{\parskip}{0pt}}
\usepackage{bookmark}
\IfFileExists{xurl.sty}{\usepackage{xurl}}{} % add URL line breaks if available
\urlstyle{same}
% fallback for those not using the hyperref driver hyperxmp:
\makeatletter
\@ifundefined{xmpquote}{\newcommand{\xmpquote}[1]{#1}}{}
\makeatother
\hypersetup{
  hidelinks,
  pdfcreator={LaTeX via pandoc}}

\author{}
\date{}

\begin{document}

\section{Bab 6: Audit Forensik Metodologi
D3TLH}\label{bab-6-audit-forensik-metodologi-d3tlh}

\textbf{CELIOS --- Center of Economic and Law Studies}

\emph{Fase 1: Evaluasi Kebijakan Ekstraktif - Pembuktian Terbalik}

\begin{center}\rule{0.5\linewidth}{0.5pt}\end{center}

\begin{quote}
\textbf{Kesimpulan Eksekutif}

D3TLH dan AMDAL telah gagal dan mati sebagai alat pelindung nyawa ruang
hidup. Dokumen-dokumen perizinan tersebut telah mereduksi penderitaan
manusia menjadi sekadar angka-angka spasial di atas kertas, berfungsi
tak lebih dari ``stempel birokrasi'' untuk melegalkan pengrusakan
ekologis secara sistematis.
\end{quote}

\begin{center}\rule{0.5\linewidth}{0.5pt}\end{center}

\subsection{Ringkasan Audit Kritis
D3TLH}\label{ringkasan-audit-kritis-d3tlh}

{\def\LTcaptype{none} % do not increment counter
\begin{longtable}[]{@{}
  >{\raggedright\arraybackslash}p{(\linewidth - 6\tabcolsep) * \real{0.2500}}
  >{\raggedright\arraybackslash}p{(\linewidth - 6\tabcolsep) * \real{0.2500}}
  >{\raggedright\arraybackslash}p{(\linewidth - 6\tabcolsep) * \real{0.2500}}
  >{\raggedright\arraybackslash}p{(\linewidth - 6\tabcolsep) * \real{0.2500}}@{}}
\toprule\noalign{}
\begin{minipage}[b]{\linewidth}\raggedright
Dimensi Audit
\end{minipage} & \begin{minipage}[b]{\linewidth}\raggedright
Skor Kerusakan
\end{minipage} & \begin{minipage}[b]{\linewidth}\raggedright
Status Vonis
\end{minipage} & \begin{minipage}[b]{\linewidth}\raggedright
Indikator Utama
\end{minipage} \\
\midrule\noalign{}
\endhead
\bottomrule\noalign{}
\endlastfoot
\textbf{Daya Tampung Udara} & \textbf{9.1 / 10} & \textbf{DAYA TAMPUNG
MELAMPAUI} & Kegagalan Deteksi Morbiditas Akumulatif \emph{(Kapasitas
PLTU: 9,825 MW / NO2 NASA: 5.76e-06 / Rasio ISPA: 2.0x)} \\
\textbf{Daya Tampung Air} & \textbf{7.4 / 10} & \textbf{DAYA TAMPUNG
MELAMPAUI} & Kegagalan Pengukuran Toksisitas \emph{(IKA Sulteng: 62.1 /
Kasus Diare: 814,407 / Konflik Air: 82)} \\
\textbf{Daya Dukung Lahan} & \textbf{10.0 / 10} & \textbf{KRISIS RUANG
DARAT} & Kegagalan Mengukur Efek Domino Lanskap \emph{(Bencana: 1,557
Kejadian / Deforestasi: 1,148,635 Ha)} \\
\textbf{Daya Dukung Sosial} & \textbf{8.7 / 10} & \textbf{DARURAT
AGRARIA} & Ilusi Jasa Budaya \& Kedaulatan Ruang \emph{(Konflik Lahan:
565 Kasus TanahKita)} \\
\textbf{Veto Kebijakan} & \textbf{10.0 / 10} & \textbf{REGULATORY
CAPTURE} & Kegagalan Tata Kelola Negara \emph{(574 Izin Baru \& 12.2 GW
PLTU Captive Diloloskan)} \\
\end{longtable}
}

\begin{center}\rule{0.5\linewidth}{0.5pt}\end{center}

\subsection{1. Filosofi Audit Forensik (Sistem Alarm
Rakyat)}\label{filosofi-audit-forensik-sistem-alarm-rakyat}

AMDAL dan D3TLH pemerintah mengklaim bersifat ``prediktif''---menilai
batasan daya dukung alam \emph{sebelum} izin diberikan. Namun, data
lapangan membuktikan bahwa dokumen-dokumen tersebut secara sistematis
cacat dan gagal melindungi ruang hidup masyarakat.

\textbf{Standpoint Riset ECC:} Kita melakukan \textbf{Pembuktian
Terbalik}. Kita tidak perlu berdebat soal rumus ``daya dukung spasial''
milik konsultan korporasi. Fakta empiris bahwa \textbf{kasus ISPA
meroket, banjir bandang rutin terjadi, konflik berdarah bereskalasi, dan
izin terus diobral secara anomali} adalah \textbf{Bukti Mutlak
(Definitive Proof)} bahwa daya dukung ekologis dan sosial wilayah
tersebut \textbf{SUDAH MELAMPAUI}.

Halaman ini merangkum semua indikator krisis menjadi sebuah palu godam
untuk memvonis bahwa sistem AMDAL/D3TLH saat ini sekadar ``stempel
birokrasi'' yang buta terhadap penderitaan manusia.

\begin{center}\rule{0.5\linewidth}{0.5pt}\end{center}

\subsection{2. Fakta: Metodologi Resmi D3TLH Pemerintah (Jasa
Ekosistem)}\label{fakta-metodologi-resmi-d3tlh-pemerintah-jasa-ekosistem}

Berdasarkan dokumen pedoman teknis D3TLH (seperti Permen LH 17/2009 dan
panduan KLHK), pemerintah saat ini menyusun D3TLH dengan pendekatan
murni spasial/bio-fisik yang disebut \textbf{Jasa Ekosistem (Ecosystem
Services)}.

Indikator resmi yang digunakan dibagi menjadi 4 kategori: * \textbf{Jasa
Penyediaan (Provisioning):} Kapasitas lahan menyediakan pangan, air
bersih, dll. * \textbf{Jasa Pengaturan (Regulating):} Kapasitas tata
air, mitigasi iklim, mitigasi banjir, pemurnian udara. * \textbf{Jasa
Pendukung (Supporting):} Siklus hara, pembentukan tanah. * \textbf{Jasa
Budaya (Cultural):} Estetika alam, rekreasi.

\subsubsection{Letak Cacat Metodologi (Blind
Spots):}\label{letak-cacat-metodologi-blind-spots}

Rumus utama yang dipakai pemerintah untuk menghitung indeks di atas
hanyalah: \textbf{Peta Ekoregion + Peta Tutupan Lahan (Land Cover)}.

\begin{itemize}
\tightlist
\item
  \textbf{Abaikan Nyawa \& Morbiditas:} Menghitung kapasitas udara dari
  peta vegetasi, namun \textbf{TIDAK PERNAH} menghitung rekam medis
  warga (ISPA) yang paru-parunya rusak akibat debu smelter.
\item
  \textbf{Abaikan Kedaulatan Ruang:} Mengukur kapasitas pertanian, tapi
  abai terhadap perampasan lahan yang memicu konflik sosial berdarah.
\item
  \textbf{Bukan Veto Kebijakan:} Saat D3TLH menyatakan daya dukung
  turun, instrumen ini tidak dipakai untuk ``menyetop'' penerbitan IUP
  (Izin Usaha Pertambangan) baru.
\end{itemize}

\begin{center}\rule{0.5\linewidth}{0.5pt}\end{center}

\subsection{3. Matriks Pembuktian Terbalik: D3TLH vs Fakta
Lapangan}\label{matriks-pembuktian-terbalik-d3tlh-vs-fakta-lapangan}

Di sinilah seluruh temuan riset kita diintegrasikan untuk
``menelanjangi'' cacat bawaan D3TLH. Di bawah ini adalah benturan
langsung antara \textbf{Mitos (Klaim Dokumen Resmi)} versus
\textbf{Realitas Lapangan (Bukti Forensik)}.

\begin{center}\rule{0.5\linewidth}{0.5pt}\end{center}

\subsubsection{A. Audit D3TLH: Daya Tampung
Udara}\label{a.-audit-d3tlh-daya-tampung-udara}

\begin{quote}
\textbf{Klaim Mitos:} ``Daya tampung udara (berdasarkan peta tutupan
lahan) diklaim masih luas dan mampu menyerap emisi.''

\textbf{Fakta Forensik ECC:} Lonjakan drastis persentase Kasus ISPA dan
penyakit saluran pernapasan di lingkar tambang. \textbf{Akumulasi Skor
Kerusakan:} \textbf{9.1 / 10} --- \emph{STATUS: DAYA TAMPUNG MELAMPAUI}
\textbf{VONIS:} Kegagalan Deteksi Morbiditas Akumulatif
\end{quote}

\paragraph{1. Korelasi PLTU \& Kualitas
Udara}\label{korelasi-pltu-kualitas-udara}

Pemerintah sebelumnya mengklaim IKU `Masih Aman'. Namun pantauan
independen Satelit TROPOMI NASA mengungkap realitas lain: konsentrasi
gas beracun NO2 meledak meroket sejajar dengan ekspansi PLTU captive.
\textbf{Threshold Kritis NASA: NO2 \textgreater{} 6.0e-6 mol/m²}.

\begin{figure}
\centering
\pandocbounded{\includegraphics[keepaspectratio,alt={PLTU vs NO2 NASA}]{visuals_bab6/chart_6_1a_pltu_no2.png}}
\caption{PLTU vs NO2 NASA}
\end{figure}

\paragraph{2. Dampak Kasus
ISPA/Pneumonia}\label{dampak-kasus-ispapneumonia}

Dokumen daya dukung mengabaikan lonjakan tajam pasien ISPA di RSUD
Morowali dan Kendari. Grafik membuktikan bahwa tren ISPA di provinsi
non-tambang relatif stabil, namun meroket secara paralel dengan asap di
provinsi sentra nikel.

\begin{figure}
\centering
\pandocbounded{\includegraphics[keepaspectratio,alt={Insiden ISPA per 10.000 Penduduk}]{visuals_bab6/chart_6_1b_ispa_trend.png}}
\caption{Insiden ISPA per 10.000 Penduduk}
\end{figure}

\paragraph{3. Fakta Beban Limbah \&
Emisi}\label{fakta-beban-limbah-emisi}

Data perizinan D3TLH fokus pada syarat emisi cerobong di atas kertas,
tetapi mengabaikan gunung-gunung debu slag (fly ash) di darat yang bebas
tertiup angin memapari puluhan desa setiap harinya. \textbf{Threshold
Kritis: 30 Juta Ton/Tahun} = 7\% dari total neraca B3 nasional 427 juta
ton dari 1 provinsi (anomali 2,4x proporsional).

\begin{figure}
\centering
\pandocbounded{\includegraphics[keepaspectratio,alt={Beban Timbulan B3 per Provinsi}]{visuals_bab6/chart_6_1c_b3_beban.png}}
\caption{Beban Timbulan B3 per Provinsi}
\end{figure}

\paragraph{4. Hilangnya Paru-Paru Udara (Emisi
CO2)}\label{hilangnya-paru-paru-udara-emisi-co2}

Audit resmi pemerintah hanya menghitung `emisi yang keluar dari corong
pabrik', tetapi dengan sengaja mengaburkan `emisi dari jutaan pohon yang
mati' akibat ekspansi lahan tambang itu sendiri. \textbf{Threshold
Kritis: 150 Juta Ton CO2e} = melampaui target NDC FOLU Net Sink 2030
(-140 juta ton CO2e).

\begin{figure}
\centering
\pandocbounded{\includegraphics[keepaspectratio,alt={Emisi CO2 Akibat Deforestasi}]{visuals_bab6/chart_6_1d_co2_emisi.png}}
\caption{Emisi CO2 Akibat Deforestasi}
\end{figure}

\begin{center}\rule{0.5\linewidth}{0.5pt}\end{center}

\subsubsection{B. Audit D3TLH: Daya Tampung
Air}\label{b.-audit-d3tlh-daya-tampung-air}

\begin{quote}
\textbf{Klaim Mitos:} ``Pembuangan tailing diizinkan selama beban
cemaran sungai/laut masih secara teori mampu mengencerkan.''

\textbf{Fakta Forensik ECC:} Penurunan drastis Indeks Kualitas Air dan
hancurnya pesisir ditandai ledakan morbiditas air. \textbf{Akumulasi
Skor Kerusakan:} \textbf{7.4 / 10} --- \emph{STATUS: DAYA TAMPUNG
MELAMPAUI} \textbf{VONIS:} Kegagalan Pengukuran Toksisitas
\end{quote}

\paragraph{1. Kualitas Air (IKA)}\label{kualitas-air-ika}

Klaim sungai/laut mampu mengencerkan limbah berbanding terbalik dengan
hancurnya Indeks Kualitas Air BPS hingga menyentuh batas cemar kotor.

\begin{figure}
\centering
\pandocbounded{\includegraphics[keepaspectratio,alt={Indeks Kualitas Air}]{visuals_bab6/chart_6_2a_ika_line.png}}
\caption{Indeks Kualitas Air}
\end{figure}

\paragraph{2. Morbiditas Diare}\label{morbiditas-diare}

AMDAL gagal menghitung dampak kontaminasi logam berat ke air tanah yang
dikonsumsi warga, dibuktikan dengan ledakan pasien Diare di lingkar
tambang. \textbf{Threshold Kritis: Incidence Rate Ratio (IRR)
\textgreater{} 2.0}.

\begin{figure}
\centering
\pandocbounded{\includegraphics[keepaspectratio,alt={Kasus Diare}]{visuals_bab6/chart_6_2b_diare_area.png}}
\caption{Kasus Diare}
\end{figure}

\paragraph{3. Konflik Nelayan \& Pesisir}\label{konflik-nelayan-pesisir}

Ekosistem tangkap nelayan dihancurkan oleh limbah tailing dan
privatisasi pesisir untuk Smelter, memicu lonjakan konflik agraria laut.

\begin{figure}
\centering
\pandocbounded{\includegraphics[keepaspectratio,alt={Konflik Nelayan \& Pesisir}]{visuals_bab6/chart_6_2c_konflik_nelayan.png}}
\caption{Konflik Nelayan \& Pesisir}
\end{figure}

\paragraph{4. Beban Tailing (Treemap
B3)}\label{beban-tailing-treemap-b3}

Resiko kebocoran Tailings Dam (Bendungan Tailing) atau Deep Sea Tailing
Placement (DSTP) yang ditutupi oleh klaim `mitigasi teknologi'.
\textbf{Threshold Kritis: 25 Juta Ton/Tahun}.

\begin{figure}
\centering
\pandocbounded{\includegraphics[keepaspectratio,alt={Proporsi Beban Limbah Tailing \& B3}]{visuals_bab6/chart_6_2d_tailing_treemap.png}}
\caption{Proporsi Beban Limbah Tailing \& B3}
\end{figure}

\begin{center}\rule{0.5\linewidth}{0.5pt}\end{center}

\subsubsection{C. Audit D3TLH: Daya Dukung
Lahan}\label{c.-audit-d3tlh-daya-dukung-lahan}

\begin{quote}
\textbf{Klaim Mitos:} ``Daya dukung lahan dan tata air tanah dinilai
aman secara matematis karena rasio ekoregion hutan dianggap masih
mencukupi.''

\textbf{Fakta Forensik ECC:} Hancurnya sabuk hijau alam memicu rentetan
bencana hidrometeorologi parah di lingkar tambang, menabrak batas fungsi
kawasan lindung. \textbf{Skor Kerusakan Lahan:} \textbf{10.0 / 10} ---
\emph{STATUS: KRISIS RUANG DARAT} \textbf{VONIS:} Kegagalan Menjaga
Fungsi Lanskap
\end{quote}

\paragraph{1. Bencana Banjir \& Longsor
(BNPB)}\label{bencana-banjir-longsor-bnpb}

Data BNPB membuktikan bahwa klaim `mitigasi bencana' dalam AMDAL sama
sekali tidak terbukti di lapangan.

\begin{figure}
\centering
\pandocbounded{\includegraphics[keepaspectratio,alt={Bencana Hidrometeorologi}]{visuals_bab6/chart_6_3a_bencana_bnpb.png}}
\caption{Bencana Hidrometeorologi}
\end{figure}

\paragraph{2. Deforestasi Primer (GFW)}\label{deforestasi-primer-gfw}

Hutan primer yang berfungsi sebagai jasa penyediaan air dan penyerap
karbon ditebang habis atas nama IUP.

\begin{figure}
\centering
\pandocbounded{\includegraphics[keepaspectratio,alt={Laju Deforestasi Pertambangan \& Sawit}]{visuals_bab6/chart_6_3b_deforestasi_gfw.png}}
\caption{Laju Deforestasi Pertambangan \& Sawit}
\end{figure}

\paragraph{3. Pelanggaran Kawasan
Lindung}\label{pelanggaran-kawasan-lindung}

Temuan \textbf{paling mematikan}: Data GFW membuktikan bahwa
\textbf{100\% dari setiap Ha deforestasi} yang terjadi di Sulteng dan
Sultra selama 10 tahun (2014--2023) terjadi di dalam \textbf{Kawasan
Lindung / Protected Areas (IUCN)}. Tidak ada satu pun hektar yang
dibabat di luar batas kawasan yang seharusnya tidak boleh disentuh.

\begin{figure}
\centering
\pandocbounded{\includegraphics[keepaspectratio,alt={Deforestasi Kawasan Lindung}]{visuals_bab6/chart_6_3c_kawasan_lindung.png}}
\caption{Deforestasi Kawasan Lindung}
\end{figure}

\paragraph{4. Aktor Deforestasi}\label{aktor-deforestasi}

Data atribusi GFW mematahkan alibi `ladang berpindah'. Pertambangan dan
Sawit adalah aktor dominan penghancur hutan.

\begin{figure}
\centering
\pandocbounded{\includegraphics[keepaspectratio,alt={Drivers of Deforestation}]{visuals_bab6/chart_6_3d_drivers_pie.png}}
\caption{Drivers of Deforestation}
\end{figure}

\begin{center}\rule{0.5\linewidth}{0.5pt}\end{center}

\subsubsection{D. Audit D3TLH: Daya Dukung
Sosial}\label{d.-audit-d3tlh-daya-dukung-sosial}

\begin{quote}
\textbf{Klaim Mitos:} ``Status kawasan dialokasikan untuk peruntukan
tambang dengan klaim bahwa masyarakat telah memberikan persetujuan
(FPIC) dalam sosialisasi amdal.''

\textbf{Fakta Forensik ECC:} Alur penindasan terbukti jelas: Persetujuan
dimanipulasi, ruang hidup jutaan hektar dirampas, dan penolakan
dibungkam dengan bui. \textbf{Skor Kerusakan Sosial:} \textbf{8.7 / 10}
--- \emph{STATUS: KRISIS KEMANUSIAAN} \textbf{VONIS:} Ilusi Kedaulatan
Ruang Warga
\end{quote}

\paragraph{1. Manipulasi Persetujuan
FPIC}\label{manipulasi-persetujuan-fpic}

`Persetujuan Warga' hanyalah stempel karet. Data investigasi Konsorsium
Pembaruan Agraria membuktikan perusahaan memanipulasi persetujuan (FPIC)
sejak fase sosialisasi AMDAL.

\begin{figure}
\centering
\pandocbounded{\includegraphics[keepaspectratio,alt={Timeline Konflik Tambang vs Masalah Izin}]{visuals_bab6/chart_6_4a_fpic_timeline.png}}
\caption{Timeline Konflik Tambang vs Masalah Izin}
\end{figure}

\paragraph{2. Perampasan Ruang Hidup}\label{perampasan-ruang-hidup}

Setelah izin keluar lewat manipulasi, perampasan paksa terjadi. Ruang
hidup warga menyusut drastis, memicu letusan konflik yang berdampak pada
ratusan ribu korban jiwa.

\begin{figure}
\centering
\pandocbounded{\includegraphics[keepaspectratio,alt={Frekuensi Letusan Konflik Perampasan Lahan}]{visuals_bab6/chart_6_4b_perampasan_lahan.png}}
\caption{Frekuensi Letusan Konflik Perampasan Lahan}
\end{figure}

\paragraph{3. Kriminalisasi Warga}\label{kriminalisasi-warga}

Di fase akhir, ketika warga melakukan penolakan yang sah atas
perampasan, negara tidak hadir melindungi, melainkan mengirim aparat
untuk memenjarakan mereka.

\begin{figure}
\centering
\pandocbounded{\includegraphics[keepaspectratio,alt={Insiden Kriminalisasi \& Kekerasan}]{visuals_bab6/chart_6_4c_kriminalisasi.png}}
\caption{Insiden Kriminalisasi \& Kekerasan}
\end{figure}

\paragraph{4. Defisit Layanan Dasar
(Faskes)}\label{defisit-layanan-dasar-faskes}

Di tengah ekspor nikel sentra Sulawesi yang meledak ratusan kali lipat,
kualitas layanan dasar hancur. Mayoritas Puskesmas gagal memenuhi
standar minimal \textbf{Sarana, Prasarana, dan Alat Kesehatan (SPA)}.
Klaim AMDAL tentang `peningkatan kesejahteraan' adalah fiksi belaka.

\begin{figure}
\centering
\pandocbounded{\includegraphics[keepaspectratio,alt={Tren Jumlah Fisik Faskes}]{visuals_bab6/chart_6_4d_faskes_line.png}}
\caption{Tren Jumlah Fisik Faskes}
\end{figure}

\begin{center}\rule{0.5\linewidth}{0.5pt}\end{center}

\subsubsection{E. Audit D3TLH: Veto
Kebijakan}\label{e.-audit-d3tlh-veto-kebijakan}

\begin{quote}
\textbf{Klaim Mitos:} ``Penyusunan D3TLH adalah dokumen sakti (veto)
yang dapat membatasi izin eksploitasi jika daya dukung lingkungan telah
terlampaui.''

\textbf{Fakta Forensik ECC:} Negara mengalami kelumpuhan tata kelola
(Regulatory Capture). Izin diobral massal, perusahaan ilegal dibiarkan,
dan infrastruktur energi kotor diloloskan di episentrum krisis.
\textbf{Skor Kegagalan Tata Kelola:} \textbf{10.0 / 10} ---
\emph{STATUS: REGULATORY CAPTURE} \textbf{VONIS:} Kegagalan Supremasi
Hukum
\end{quote}

\paragraph{1. Obral Konsesi Legal}\label{obral-konsesi-legal}

Di tengah memuncaknya status krisis daya dukung lingkungan, pemerintah
secara paradoks justru menerbitkan ratusan izin eksploitasi tambang
(IUP) baru. Dokumen veto tidak berfungsi.

\begin{figure}
\centering
\pandocbounded{\includegraphics[keepaspectratio,alt={Lonjakan Penerbitan IUP Baru}]{visuals_bab6/chart_6_5a_iup_baru.png}}
\caption{Lonjakan Penerbitan IUP Baru}
\end{figure}

\paragraph{2. Pembiaran Pelanggaran
Korporat}\label{pembiaran-pelanggaran-korporat}

Bukti mutlak `Regulatory Capture'---bahkan ketika perusahaan beroperasi
ilegal, menabrak izin, tumpang tindih, atau HGU kedaluwarsa, negara
tidak berani melakukan penegakan hukum dan membiarkannya.

\begin{figure}
\centering
\pandocbounded{\includegraphics[keepaspectratio,alt={Distribusi Modus Pelanggaran Izin Korporat}]{visuals_bab6/chart_6_5b_modus_pelanggaran.png}}
\caption{Distribusi Modus Pelanggaran Izin Korporat}
\end{figure}

\paragraph{3. Karpet Merah Energi Kotor (PLTU
Captive)}\label{karpet-merah-energi-kotor-pltu-captive}

Inkonsistensi paling telanjang terhadap komitmen iklim. Di wilayah
ekoregion krisis, pemerintah memberikan karpet merah pembangunan
infrastruktur penyumbang emisi terbesar (PLTU Batubara Captive) khusus
untuk menyuplai kawasan smelter nikel.

\begin{figure}
\centering
\pandocbounded{\includegraphics[keepaspectratio,alt={Proporsi Status PLTU Captive}]{visuals_bab6/chart_6_5c_pltu_status_pie.png}}
\caption{Proporsi Status PLTU Captive}
\end{figure}

\end{document}
